\documentclass[12pt,a4paper,twoside]{article}
\usepackage[utf8]{inputenc}
\usepackage[T1]{fontenc}
\usepackage{amsmath, esint}
\usepackage{amsfonts}
\usepackage{listings}
\usepackage{amssymb}
\usepackage{soul, color}
\usepackage{cancel}
\usepackage{wrapfig}
\usepackage[left=2.54cm, right=2.54cm, top=2.54cm, bottom=2.54cm]{geometry}
\usepackage[colorlinks,linkcolor=cyan]{hyperref}
\linespread{1.0}
\usepackage{graphicx}
%%%%%%%%%%%%%%%%%%%%%%%%%%%%%%%%%%%%%%%%%%%%%%%%%%%%%%%%%%%%%%%%%%%%%%%%%%%%%%%%%%%%%%%%%%%%%%
% Define the header.
\usepackage{fancyhdr}
\pagestyle{fancy}
\fancyhf{}
\fancyhead[RE,RO]{\leftmark} % Define header: LE= Left in Even pages, RE = Right in Even pages. CO = Center in Odd pages. \leftmark = section number and section name
\fancyhead[LE,LO]{\LaTeX\ \textsc{Template for Notes}}
\fancyfoot[CE,CO]{\thepage}
%%%%%%%%%%%%%%%%%%%%%%%%%%%%%%%%%%%%%%%%%%%%%%%%%%%%%%%%%%%%%%%%%%%%%%%%%%%%%%%%%%%%%%%%%%%%%%
% Define Griffith's script r
\usepackage{tikz}
\usetikzlibrary{arrows,scopes}
\newcommand{\rc}{
\resizebox{!}{1.25ex}{
    \begin{tikzpicture}[>=round cap]
        \clip (0.09em,-0.05ex) rectangle (0.61em,0.81ex);
        \draw [line width=.11ex, <->, rounded corners=0.13ex] (0.1em,0.1ex) .. controls (0.24em,0.4ex) .. (0.35em,0.8ex) .. controls (0.29em,0.725ex) .. (0.25em,0.6ex) .. controls (0.7em,0.8ex) and (0.08em,-0.4ex) .. (0.55em,0.25ex);
    \end{tikzpicture}
}
}

\newcommand{\brc}{
\resizebox{!}{1.3ex}{
    \begin{tikzpicture}[>=round cap]
        \clip (0.085em,-0.1ex) rectangle (0.61em,0.875ex);
        \draw [line width=.2ex, <->, rounded corners=0.13ex] (0.1em,0.1ex) .. controls (0.24em,0.4ex) .. (0.35em,0.8ex) .. controls (0.29em,0.725ex) .. (0.25em,0.6ex) .. controls (0.7em,0.8ex) and (0.08em,-0.4ex) .. (0.55em,0.25ex);
    \end{tikzpicture}
}
}
\newcommand{\hrc}{\hat{\brc}}
%%%%%%%%%%%%%%%%%%%%%%%%%%%%%%%%%%%%%%%%%%%%%%%%%%%%%%%%%%%%%%%%%%%%%%%%%%%%%%%%%%%%%%%%%%%%%%
\numberwithin{equation}{section} % Make the equation numbering become (section number.equation #)
%%%%%%%%%%%%%%%%%%%%%%%%%%%%%%%%%%%%%%%%%%%%%%%%%%%%%%%%%%%%%%%%%%%%%%%%%%%%%%%%%%%%%%%%%%%%%%
% Code listing from Overleaf.
\definecolor{codegreen}{rgb}{0,0.6,0}
\definecolor{codegray}{rgb}{0.5,0.5,0.5}
\definecolor{codepurple}{rgb}{0.58,0,0.82}
\definecolor{backcolour}{rgb}{0.95,0.95,0.92}

\lstdefinestyle{mystyle}{
    backgroundcolor=\color{backcolour},   
    commentstyle=\color{codegreen},
    keywordstyle=\color{magenta},
    numberstyle=\tiny\color{codegray},
    stringstyle=\color{codepurple},
    basicstyle=\ttfamily\footnotesize,
    breakatwhitespace=false,         
    breaklines=true,                 
    captionpos=b,                    
    keepspaces=true,                 
    numbers=left,                    
    numbersep=5pt,                  
    showspaces=false,                
    showstringspaces=false,
    showtabs=false,                  
    tabsize=2
}
\lstset{style=mystyle}
%%%%%%%%%%%%%%%%%%%%%%%%%%%%%%%%%%%%%%%%%%%%%%%%%%%%%%%%%%%%%%%%%%%%%%%%%%%%%%%%%%%%%%%%%%%%%%

\begin{document}
\begin{center}
		\textbf{\Large{\LaTeX \ Template for Math and Physics notes}}\\
		\textit{Kuzma Cosmos (2022)}
	\end{center}

\section{Some complicated math notations}

\subsection{Jacobian}
This is the equation for changing variable in double integral:
\begin{equation}
    \iint_{\mathcal{S}}f(x,y) dxdy = \iint_{\mathcal{D}}f(x(u, v),\ g(u,v))\left|\dfrac{\partial(x,y)}{\partial(u,v)}\right|dudv
    \label{eq: change-variable-double-int}
\end{equation}
where $(x,y)=T(u,v)=(x(u,v),\ y(u,v))$ with a defined transformation $T: \mathcal{D}\to \mathcal{S}$. $\dfrac{\partial(x,y)}{\partial(u,v)}$ is called \textbf{Jacobian} of $T$, denoted by $\mathrm{Jac}(T)$:

\begin{equation}
    \boxed{
    \mathrm{Jac}(T) = \frac{\partial(x,y)}{\partial(u,v)} =
    \def\arraystretch{2}
    \begin{vmatrix}
    \dfrac{\partial x}{\partial u} & \dfrac{\partial x}{\partial v}\\ 
    \dfrac{\partial y}{\partial u} & \dfrac{\partial y}{\partial v}
    \end{vmatrix}
    }
\end{equation}

\newpage

\section{Some physics examples}
\subsection{Biot-Savart Law}
This is equation for Biot-Savart Law in magnetostatics, from Griffiths' Textbook.
\begin{equation}
    \boxed{
    \overrightarrow{B} = \dfrac{\mu_0}{4\pi}\int_{\mathcal{W}}\dfrac{\overrightarrow{I}\times \hrc}{\rc^2}dl=\dfrac{\mu_0I}{4\pi}\int_{\mathcal{W}}\dfrac{d\overrightarrow{l}\times \hrc}{\rc^2}
    }
    \label{Eq: Biot-Savart}
\end{equation}


\newpage

\section{Python codes in \LaTeX}
Sample python code.
\lstinputlisting[language = Python]{sample.py}

% Write python code on LaTeX, example excerpted from Overleaf
\begin{lstlisting}[language=Python, caption=Python example]
import numpy as np
    
def incmatrix(genl1,genl2):
    m = len(genl1)
    n = len(genl2)
    M = None #to become the incidence matrix
    VT = np.zeros((n*m,1), int)  #dummy variable
    
    #compute the bitwise xor matrix
    M1 = bitxormatrix(genl1)
    M2 = np.triu(bitxormatrix(genl2),1) 

    for i in range(m-1):
        for j in range(i+1, m):
            [r,c] = np.where(M2 == M1[i,j])
            for k in range(len(r)):
                VT[(i)*n + r[k]] = 1;
                VT[(i)*n + c[k]] = 1;
                VT[(j)*n + r[k]] = 1;
                VT[(j)*n + c[k]] = 1;
                
                if M is None:
                    M = np.copy(VT)
                else:
                    M = np.concatenate((M, VT), 1)
                
                VT = np.zeros((n*m,1), int)
    
    return M
\end{lstlisting}
\end{document}